\chapter{Родитель размерностной трансмутации такта рождения}
\label{ch:v10-nonequilibrium-bath-birth-tick-dimensional-transmutation}

Проверим самый сильный внутренний кандидат предыдущего аудита. Унаследованные
данные относительного $U(1)$-сектора имеют вид
\begin{equation}
 b=2,\qquad g^2(\mu_{\rm spec})=\frac38,
 \qquad
 \log\frac{\Lambda_L}{\mu_{\rm spec}}=\frac{32\pi^2}{3}.
\end{equation}
Если частоту рождения условно отождествить с ландауовской частотой,
\begin{equation}
 \omega_{\rm DT}=c\Lambda_L
 =c\mu_{\rm spec}\exp\!\left(\frac{32\pi^2}{3}\right),
\end{equation}
то получается относительный такт
\begin{equation}
 \tau_{\rm DT}=\frac1{\omega_{\rm DT}}
 =\frac{\exp(-32\pi^2/3)}{c\mu_{\rm spec}}.
\end{equation}
Трансмутация действительно создаёт размерный объект из безразмерной связи,
но не создаёт исходную размерную единицу: при
$\mu_{\rm spec}\mapsto\mu_{\rm spec}/s$ имеем
$\tau_{\rm DT}\mapsto s\tau_{\rm DT}$.

В нормированных логарифмических невязках существует строгий общий родитель
\begin{equation}
 \mathcal P_{\rm DT}=\frac12\left[
 \left(x_{\rm RG}-\frac{32\pi^2}{3}\right)^2
 +x_{\rm tick}^2+x_{\rm transit}^2\right],
 \qquad \nabla^2\mathcal P_{\rm DT}=I_3.
\end{equation}
Он условно связывает RG-отношение, $\tau_{\rm birth}\omega_{\rm DT}=1$ и
$\ell_{\rm edge}=c\tau_{\rm birth}$. Однако размерная карта имеет
ранг/ядро $3/2$, после фиксации $c$ --- $4/1$, и лишь независимый якорь
$\mu_{\rm spec}$ даёт $5/0$.

Более того, наивная типизация $\mu_{\rm spec}=\Lambda_{43}$ противоречит
уже закрытой геометрической связи. Трансмутация требует
\begin{equation}
 \tau_{\rm birth}c\Lambda_{43}=\exp\!\left(-\frac{32\pi^2}{3}\right),
\end{equation}
тогда как K43-ветвь даёт точно
\begin{equation}
 \tau_{\rm birth}c\Lambda_{43}=42.
\end{equation}
Их отношение равно $42\exp(32\pi^2/3)\ne1$.

\begin{theorem}[Относительная, но не абсолютная трансмутация такта]
Унаследованная RG-пара $(b,g^2)=(2,3/8)$ однозначно задаёт отношение
$\tau_{\rm DT}c\mu_{\rm spec}=\exp(-32\pi^2/3)$, но сохраняет одну
масштабную моду после фиксации $c$. Поэтому абсолютный такт рождения из неё
не следует.
\end{theorem}

\begin{nogocriterion}[Нельзя одновременно отождествить $\mu_{\rm spec}$ с K43-обрезанием]
Без нового граничного условия равенства
$\tau c\mu_{\rm spec}=\exp(-32\pi^2/3)$ и
$\tau c\Lambda_{43}=42$ несовместимы при
$\mu_{\rm spec}=\Lambda_{43}$. Такое отождествление не является физической
типизацией; оно создаёт точное противоречие.
\end{nogocriterion}

\begin{protocol}\begin{enumerate}
\item условная архитектура: \textbf{$9/9$};
\item унаследованные RG-данные: \textbf{$2/2$};
\item типизация RG-связи в носитель ванны: \textbf{$0/1$};
\item размерная карта после фиксации $c$: \textbf{ранг/ядро $4/1$};
\item совместимость наивного отождествления с K43: \textbf{$0/1$};
\item физическое происхождение пакета: \textbf{$0/4$};
\item происхождение абсолютного такта: \textbf{$0/1$};
\item следующий гейт: \textbf{K43--RG согласование граничного условия}.
\end{enumerate}\end{protocol}

Машинный сертификат сохранён в
\path{s2t_v10_cell_birth_four_volume_nonequilibrium_bath_birth_tick_dimensional_transmutation_parent_origin_gate_results.json}.